% Philippica XII (philippic-12) --- English translation
% Translated by Claude (Anthropic) from the Latin (Perseus).
% Style guide: STYLE.md.

\ciceroSection{1} Although it seems least becoming, senators, that a man to whom you so often assent on matters of the greatest moment should be fooled, deceived, led into error, I console myself nonetheless, since I went astray together with you, side by side, and along with our most wise consul. For when two men of consular rank had brought us the hope of an honourable peace, it seemed --- because they were close friends of Marcus Antonius, members of his household --- that they knew some wound of his that was unknown to us. In the one's house were a wife and children; the other was forever sending letters and receiving them, and openly taking Antony's side.

\ciceroSection{2} These men, suddenly urging peace --- a thing they had not done for a long time --- seemed to do so not without reason. And a consul joined them as their advocate. And what a consul! If we look for prudence, one least capable of being deceived; if for valour, one who would approve no peace save with a foe yielding and beaten; if for greatness of spirit, one who would set death before slavery. But you, senators, seemed not so much forgetful of your own most weighty decrees as, when the hope of a surrender --- which his friends would rather call a peace --- had been held out, to be thinking of terms to be imposed, not accepted. And my own hope had been heightened --- yours too, I suppose --- because I kept hearing that Antony's house was stricken with grief, that his wife was lamenting. Here, too, I saw Antony's supporters --- on whose faces my eyes dwell --- looking sadder than usual.

\ciceroSection{3} But if it is not so, why was mention of peace made by Piso and Calenus of all men, why at this time, why so unforeseen, why so suddenly? Piso denies that he knows anything, denies that he has heard a word; Calenus denies that any fresh news has been brought. And they deny it now --- now that they think they have entangled us in a peace-making embassy. What need, then, of a new plan, if in the matter there is nothing new at all? We have been deceived, I say, senators: it was Antony's cause that his friends pleaded, not the public's. This indeed I saw, but as if through a mist: the safety of Decimus Brutus had dazzled the edge of my mind. And if substitutes were wont to be given in war, I would gladly suffer myself to be shut up in his place, so that Decimus Brutus might be let out. And by this saying of Quintus Fufius we were caught:

\ciceroSection{4} `Shall we not hear Antony even if he withdraws from Mutina --- not even if he says that he will put himself in the Senate's power?' It seemed harsh: and so we were broken, we gave way. Does he, then, withdraw from Mutina? `I do not know.' Does he obey the Senate? `I believe so,' says Calenus, `but in such a way as to keep his dignity.' By Hercules, you must labour mightily, senators, to lose your own dignity, which is the greatest there is, and to keep Antony's --- which neither exists nor can exist --- so that he may recover through you the dignity he has lost through himself. If he were treating with you from a prostrate position, I might perhaps give ear --- and yet --- but I would rather put it this way: I would give ear. A man on his feet must be resisted, or our liberty surrendered along with our dignity. `But the matter is no longer open: the embassy has been resolved upon.'

\ciceroSection{5} But what is not still open to a wise man, if it can be put right? Any man may err; none but a fool persists in his error. For second thoughts, as they say, are wont to be the wiser. That mist I spoke of a moment ago has been scattered; the light has broken, all lies open, we see everything --- and not by ourselves alone, but we are warned by our own friends. You marked, a little while ago, what the most excellent man's speech was. `I found my house in mourning,' he said, `my wife, my children. Good men wondered, my friends reproached me, that in hope of peace I had taken on the embassy.' And no wonder, Publius Servilius: for by your most severe and most weighty motions Antony has been stripped of all --- I will not say dignity, but even hope of safety.

\ciceroSection{6} That you should go as an envoy to him --- who would not wonder at it? I learn it from my own case: I feel how I --- whose counsel is the same as yours --- am being faulted. Are we the only ones being blamed? What? Did that most valiant man Pansa speak just now, so carefully and so long, without cause? What did he accomplish, but to drive off from himself a false suspicion of treachery? And whence comes that suspicion? From his sudden advocacy of peace, which he took up on a sudden, caught by the same error as ourselves.

\ciceroSection{7} But if an error has been made, senators, through a false and deceiving hope, let us return to the road. The best harbour for a man who repents is a change of plan. For what good, immortal gods, can our embassy do the commonwealth? Do good, do I say? What if it will even do harm? Do harm? What if it has already injured and harmed? Or do you suppose that the keenest and bravest longing of the Roman people to recover its liberty has not been lessened and weakened, now that an embassy of peace has been heard of? What do you think of the towns? What of the colonies? What of all Italy? Will it be of the same zeal with which it had blazed up against the common conflagration? Or do we not think it will come to pass that those will repent of having declared and paraded their hatred of Antony --- those who promised money, who promised arms, who gave themselves wholly, in spirit and in body, to the safety of the commonwealth? How will Capua approve this plan of ours --- Capua, which in these times is a second Rome? She judged her impious citizens, cast them out, shut them out. From that city, I say, striving most bravely, Antony was snatched out of her very hands.

\ciceroSection{8} What? Do we not, by such plans, cut the sinews of our legions? For who will be of a spirit kindled for war, when a hope of peace is held out? That very Martian legion, heavenly and divine, will at this news grow faint and soften, and will lose that most beautiful name: their swords will fall from them, their arms slip from their hands. For having followed the Senate, it will not think it ought to be more grievous in its hatred of Antony than the Senate itself. I am ashamed for this legion, ashamed for the Fourth, which with equal valour, approving our authority, abandoned Antony --- not as its consul and commander, but as an enemy and assailant of the fatherland; ashamed for that excellent army formed from the union of two, now purified by sacrifice, now marched out to Mutina; which, if it hears the name of peace --- that is, of our fear --- even if it does not draw back its foot, will surely halt.

\ciceroSection{9} For why should it hasten to fight, when the Senate calls it back and sounds the retreat? And what more unjust than that we should decree about peace without the knowledge of those who are waging the war --- and not only without their knowledge, but against their will? Or do you think that Aulus Hirtius, that most renowned consul, and Gaius Caesar, born by the gods' favour for these times --- whose letters declaring the hope of victory I hold in my hand --- wish for peace? They long to conquer; and the sweetest and most beautiful name of peace they have desired to win not by a compact, but by victory. What? With what feeling, pray, do you think Gaul will hear of this? For she holds the leading part in driving back, in conducting, in sustaining this war. Gaul, following the very nod --- not to say the command --- of Decimus Brutus, has made firm the beginnings of war with arms, with men, with money; the same Gaul has thrown its whole body in the way of Marcus Antonius's cruelty; she is drained dry, laid waste, burned: she bears all the wrongs of war with composure, so long as she may beat back the danger of slavery.

\ciceroSection{10} And --- to pass over the rest of Gaul, for all its parts are alike --- the people of Patavium shut out some of Antony's agents, drove out others, and aided our commanders with money, with soldiers, and --- what was most lacking --- with arms. The rest did the same --- those who were once in the same cause, and were thought to be estranged from the Senate on account of the wrongs of many years: and it is no wonder that these are faithful now that the commonwealth has been shared with them --- men who, even when they had no part in it, always kept their faith. To all these, then, who are hoping for victory, shall we bring the name of peace --- that is, despair of victory?

\ciceroSection{11} What if there cannot even be any peace at all? For what terms of peace are there, in which nothing can be granted to the man with whom you would make peace? By many things Antony was invited by us to peace: yet he preferred war. Envoys were sent, against my opposition --- but sent nonetheless; instructions were delivered: he did not obey. He was warned not to besiege Brutus, to withdraw from Mutina: he assailed the city all the more fiercely. And shall we send envoys of peace to a man who rejected the messengers of peace? Do we think he will be more modest in his demands face to face than he was when he sent his instructions to the Senate? And yet then he asked for things which seemed altogether shameless, but which still could in some fashion be granted; he had not yet been cut to pieces by your judgments, so weighty and so many, and by your marks of disgrace: now he asks for things we can in no way grant, unless we are first willing to confess ourselves beaten in the war.

\ciceroSection{12} We judged the decrees of the Senate that he produced to be forgeries: can we now judge them genuine? We declared that the laws were carried by violence and against the auspices, and that neither people nor plebs is bound by them: do you think they can now be restored? You judged that Antony had embezzled seven hundred million sesterces of the public money: can the peculation be free of guilt? Exemptions were sold by him to states, and priesthoods, and kingdoms: shall those tablets be fixed up again which you by your decrees took down? But if we wish to bury what we have decreed, can we also blot out the memory of these things? For when will any after-age forget through whose crime we have gone in this foul attire? Granted that the blood of the centurions of the Martian legion, spilt at Brundisium, be washed away --- can the proclamation of his cruelty be washed out? To pass over what lies between: what length of time will remove the foul monuments of his works around Mutina, the proofs of his crime, the traces of his brigandage?

\ciceroSection{13} To this savage and unclean parricide, then, what have we, immortal gods, to remit? Farther Gaul, perhaps, and an army? What else is this but to make, not peace, but a postponement of war --- nay, not only to prolong the war but even to surrender the victory? Will he not have conquered, if on any terms whatever he comes into this city with his followers? Now we hold all things by arms; in authority we are strongest; so many ruined citizens are away, who followed their wicked leader; yet the faces and the talk of those of that number who have been left behind in the city we cannot bear. What do you think will happen, when so many of them burst in at once, and we have laid down our arms while they have not laid down theirs --- shall we not, by our own counsels, be conquered for ever?

\ciceroSection{14} Set before your eyes Marcus Antonius the consular; add Lucius, hoping for the consulship; fill out the rest, who plan offices and commands --- and not of our order only: do not despise even such creatures as Numisius Tiro and Mustela Seius. A peace made with these men will not be peace, but a compact of slavery. The noble saying of Lucius Piso, that most distinguished man, was rightly praised by you, Pansa, not only in this order but also before the assembly. He said that he would depart from Italy, would leave his household gods and the seats of his fathers, if --- which heaven avert the omen! --- Antony should crush the commonwealth.

\ciceroSection{15} I ask you then, Lucius Piso: would you not think the commonwealth crushed, if men so many, so impious, so bold, so steeped in crime, were taken back? Those whom, not yet stained with such great crimes, we could scarcely endure --- do you reckon that these men, now covered over with every crime, will be tolerable to the state? Either we must take that counsel of yours, believe me --- to give way, to depart, to follow out a life of want and wandering --- or our necks must be given to the brigands, and we must fall in our own country. Where are they now, Gaius Pansa, those most beautiful exhortations of yours, by which the Senate, roused, and the Roman people, set aflame, not only heard but learned that to a Roman nothing is fouler than slavery?

\ciceroSection{16} Was it for this that we put on the soldier's cloak, took up arms, shook out all the youth from the whole of Italy --- that, with an army most flourishing and most great, envoys should be sent to sue for peace? If peace is to be accepted, why are we not asked? If it is to be demanded, what do we fear? Am I to be in this embassy, or be mixed up in that counsel, in which, even if I dissent from the rest, the Roman people will not so much as know it? So it will come about that, if anything is remitted or granted, Antony will always sin at my peril --- since the power of sinning will seem to have been granted him by me.

\ciceroSection{17} But if account had to be taken of peace with Marcus Antonius's band of brigands, still my person was the last that should have been chosen to procure that peace. I never voted that envoys should be sent; I, before the envoys' return, dared to say that, even if they should bring peace itself, it must be rejected, since under the name of peace a war lay hidden; I was the first to put on the cloak; I always called him an enemy, when others called him an adversary; always called this a war, when others called it a tumult. And this not in the Senate only: I urged the same always before the people; nor against him alone, but against the partners and the agents of his crimes, both those present and those who are with him --- in short, against the whole house of Marcus

\ciceroSection{18} Antonius I have always inveighed. And so, just as the impious citizens, eager and glad at the hope of peace held out, were congratulating one another as if they had won, so they renounced me as biased, and complained of me; they distrusted Servilius too: they remembered that Antony had been transfixed by his motions; Lucius Caesar, a brave and steadfast senator indeed --- but still an uncle; Calenus an agent; Piso an intimate; and you yourself, Pansa, that most vehement and most valiant consul, they now think have been made gentler: not that it is so, or could be, but the mention of peace made by you has brought to many the suspicion of a changed purpose. That I have been thrust in among these persons, Antony's friends take ill: and we must humour them, since once for all we have begun to be generous.

\ciceroSection{19} Let the envoys set out with the best of omens --- but let those set out at whom Antony will take no offence. But if you do not trouble yourselves about Antony, you ought at least, senators, to take thought for me. Spare my eyes at least, and grant some indulgence to a just grief. For with what face shall I be able to look upon --- I leave aside the enemy of my country, in which my hatred of him is shared with you --- but how shall I look upon an enemy most cruel to me alone, as his most bitter harangues about me declare? Do you think me so made of iron that I could meet him, or look upon a man who lately, when in a public meeting he was making gifts to those who seemed to him the boldest among the parricides, said that he was giving my property to Petusius of Urbinum --- a man flung up out of the shipwreck of a handsome patrimony onto these Antonian rocks?

\ciceroSection{20} Or shall I be able to look upon Lucius Antonius, whose cruelty I could not have escaped, had I not defended myself with the walls and gates and the zeal of my own town? And this same man --- the Asiatic murmillo, the brigand of Italy, the colleague of Lento and Nucula --- when he was giving gold coins to Aquila, a chief centurion, said that he gave them out of my goods: for if he had said out of his own, he thought not even the Eagle herself would have believed him. My eyes will not endure, I say, the sight of Saxa, of Cafo, nor the two praetors, nor the two tribunes-designate, nor Bestia, nor Trebellius, nor Titus Plancus. I cannot with composure look upon so many enemies, so savage, so steeped in crime; and this comes not from any fastidiousness of mine, but from love of the commonwealth.

\ciceroSection{21} But I will conquer my feelings and command myself: a most just grief, if I cannot break it, I will hide. What? Do you think, senators, that I ought to take some account of my life? --- which to me, indeed, is far from dear, especially now that Dolabella has made death a thing to be wished for, provided only it be without torture and torment; yet to you and to the Roman people my breath ought not to be cheap. For I am the man --- unless perhaps I deceive myself --- who by my vigils, my cares, my motions, and the perils too, very many, which I underwent because of the bitterest hatred of all the impious against me, have brought it about that I am no hindrance to the commonwealth --- to say nothing more arrogant.

\ciceroSection{22} This being so, do you think I ought to take no thought at all for my own danger? Here, when I was in the city and at home, many attempts were nonetheless often made --- here, where not only the faithfulness of my friends but the eyes of the whole state keep watch over me: what do you think, when I have set out on a journey --- a long one, at that --- will there be no ambushes to dread? There are three roads to Mutina --- whither my heart hastens, that I may as soon as possible look upon Decimus Brutus, that pledge of the Roman people's liberty; in whose embrace I would gladly breathe out the last breath of my life, now that all the actions of these months, all my motions, have arrived at the end I set before myself. Three roads, then, as I said: from the upper sea the Flaminian, from the lower the Aurelian, between them the Cassian.

\ciceroSection{23} Now, I beg you, consider whether my suspicion of danger strays from a fair guess. The Cassian Way divides Etruria. Do we know, then, Pansa, in what regions the septemviral authority of Caesennius Lento now is? With us he certainly is not, neither in mind nor in body. But if he is either at home or not far from home, he is certainly in Etruria --- that is, on my road. Who, then, will warrant to me that Lento is content with a single head? Tell me besides, Pansa, where Ventidius is --- to whom I was always a friend, before he became so openly the enemy of the commonwealth and of all good men. I can avoid the Cassian and keep to the Flaminian: but what if Ventidius has come to Ancona, as is said --- shall I be able to reach Ariminum in safety? There remains the Aurelian. Here, indeed, I even have garrisons: for there are the estates of Publius Clodius. The whole household will come out to meet me; it will invite me to its hospitality --- on account of our most well-known friendship.

\ciceroSection{24} Shall I entrust myself to these roads --- I who lately, at the Terminalia, did not dare to go out into the suburbs to return the same day? Within my own walls I scarcely keep myself safe, without the guard of my friends. And so I stay in the city; if I am allowed, I will stay. This is my station, this my watch, this my guard-post, this my standing garrison. Let others hold the camps and conduct the business of war; let them slay the enemy --- for this is the chief thing; we, as we say and have always done, will guard the city and the affairs of the city, side by side with you. Nor in truth do I refuse this duty --- though I see the Roman people refuse it on my behalf. No one is less timid than I, yet no one more cautious. The facts declare it. It is now the twentieth year that all the wicked aim at me alone. And so they have paid the penalty --- not to me, let me say, but to the commonwealth: me the commonwealth has so far kept safe for itself. I will say this with diffidence --- for I know that anything can befall a man --- yet once, when I was beset by the chosen strength of the most powerful men, I fell, and fell knowingly, that I might rise again most honourably.

\ciceroSection{25} Can I, then, seem cautious enough, foreseeing enough, if I commit myself to a journey so hostile and so perilous? Those who are engaged in public life ought, at their death, to leave behind them glory --- not the reproach of fault and the censure of folly. What good man does not mourn the death of Trebonius? Who does not grieve at the destruction of such a citizen and such a man? But there are those who say --- harshly, to be sure, yet they say it --- that there is the less cause to grieve, because he did not guard himself against an unclean and wicked man. For the man who professes himself the guardian of many, the wise say, ought first to be the guardian of his own life. When you are hedged about by the laws and by the fear of the courts, not everything is to be feared, nor need guards be sought against every ambush. For who would dare to attack a man in daylight, who on a soldiers' highroad, who when he is well attended, who when he is a man of mark? These considerations hold good neither at this time nor in my case.

\ciceroSection{26} For the man who lays violent hands on me will not only not dread punishment, but will even hope for glory and rewards from the bands of brigands. These things I can foresee in the city: it is easy to look about me --- whence I go out, whither I advance, what is on the right, what on the left. Shall I be able to do the same on the tracks of the Apennines? On which, even if there are no ambushes --- though there could most easily be --- my mind will yet be so anxious that it can attend to nothing of the embassy's duties. But suppose I have escaped the ambushes, broken through the Apennines: still I must come to a meeting and a parley with Antony. What place, pray, will be chosen? If outside the camp, let the others look to it: I think I shall scarcely be safe. I know the man's frenzy, I know his unbridled violence; whose bitterness of character and monstrousness of nature not even an admixture of wine is wont to temper. This man, inflamed with wrath and madness, with his brother Lucius at his side --- that most loathsome beast --- will surely never keep his sacrilegious and impious hands off me.

\ciceroSection{27} I remember parleys both with the bitterest enemies and with citizens most gravely at odds. Gnaeus Pompeius, son of Sextus, when consul --- I being present, for I was a recruit in his army --- parleyed with Publius Vettius Scato, the leader of the Marsi, between the two camps; and on that very day I remember that Sextus Pompeius, the consul's brother, a learned and wise man, came from Rome to the very parley. When Scato had greeted him, `What shall I call you?' he said. And the other: `By goodwill a guest, by necessity a foe.' In that parley there was fairness; no fear, no suspicion lurked beneath; the hatred, too, was moderate. For the allies were not seeking to wrest the citizenship from us, but to be received into it. Sulla, with Scipio, between Cales and Teanum --- the one having brought the flower of the nobility, the other the allies in the war --- conferred together over terms and conditions touching the authority of the Senate, the votes of the people, the rights of citizenship. That parley did not, in the end, keep faith: yet it was free from violence and danger. Can we, then, be equally safe amid Antony's brigandage? We cannot; or, if the rest can, I have no confidence that I can.

\ciceroSection{28} But if we are not to meet outside the camp, what camp shall be taken for the parley? Into ours he will never come; much less shall we come into his. It remains, then, that the demands be both received and sent back by letter. So we shall be in the camp, and my opinion on all his demands will be one and the same; and when I have stated it here, in your hearing, count me to have gone and returned: I shall have accomplished the embassy. Everything, by my motion, I will refer back to the Senate, whatever Antony shall demand. For it is not lawful otherwise, nor has it been allowed us by this order --- as is customary, by ancestral usage, when wars are finished and the matter is committed to ten commissioners --- nor have we received any instructions at all from the Senate. And when I press these things in the council, with some, as I suppose, resisting, is it not to be feared that the untrained crowd of soldiers will think that peace is being held back through me?

\ciceroSection{29} Grant that the new legions do not disapprove this counsel of mine --- for the Martian and the Fourth, which think of nothing but dignity and honour, will approve it, I know for certain: but what of the veterans? Are we not anxious how they may take my severity (for not even they themselves wish to be feared)? For they have heard many false things about me; many things have wicked men carried to them --- though their interests, as you are the best of witnesses, I have always upheld by my motions, my authority, my speech: yet they believe the wicked, they believe the turbulent, they believe their own. Brave they are, to be sure; but, on account of the memory of what they did for the liberty of the Roman people and the safety of the commonwealth, they are too fierce, and refer all our counsels to their own strength.

\ciceroSection{30} Their reflection I do not fear; their onset I dread. And if I escape these great dangers too, do you think my return will be safe enough? For when I have defended your authority in my wonted way, and shown my loyalty and constancy to the commonwealth, then I shall have to dread not only those who hate me, but those too who envy me. Let my life, then, be kept for the commonwealth, so long as dignity or nature shall allow; let it be reserved for the fatherland; let my death have either the necessity of fate, or, if it must be met before that, let it be met with glory. This being so --- although the commonwealth, to put it most mildly, has no need of this embassy --- yet if it shall be safe to go, I will set out. In a word, senators, I will measure the whole plan of this matter not by my own danger, but by the advantage of the commonwealth; and since I have free time for it, I think it should be weighed again and again, and that above all should be done which I shall judge to be most in the interest of the commonwealth.
